\section{Repeated natural heads and the smaller pair counts}
\label{sec:small-pair-heads}
For a finite group $J$, write $d_p(J)=\dim_{\F_p}J/(J'J^p)$.
All head modules below come from an actual subgroup of the same complete
$J\leq S_b$.

\begin{proposition}[Repeated natural heads with central binary lifts]
\label{prop:central-repeat}
Let
\[
 R_d^A=\F_4^d:C_3,\qquad
 R_d^S=(\F_2^2)^d:S_3,
\]
where the same scalar, respectively natural, action is used on all $d$
copies. Suppose $Q$ has a specified central elementary binary subgroup
$C$ of rank $c$ and $Q/C\cong R_d^A$ or $R_d^S$, with $d\geq1$.
For every complete $J\leq S_b$,
\[
 |\Epi(J,Q)|\leq B_T(d)2^{\max(c,d)b/2},
 \qquad B_A(d)=2\cdot4^d,\quad B_S(d)=6\cdot4^d.
\]
No factor $|\Aut Q|$ is required in this bound.
\end{proposition}
\begin{proof}
Put $r=d_2(J)$. In the $A$ case let $F=J'$, and in the $S$ case let
$F=J''$. In either case set $B=J/F$ and $D=H_1(F,\F_2)$, with its
honest $B$ action. For every onto top map $\theta$ to $C_3$ or $S_3$,
let $V_\theta$ be the natural two-dimensional binary module. In the
first case the odd part of the abelian $B$ maps onto $C_3$; in the
second case the odd part of the abelian $B'$ does so. This is a normal
odd subgroup with zero invariants on $V_\theta$. Averaging and
inflation--restriction therefore give
\[
 H^i(B,V_\theta)=0\quad(i\geq0),\qquad
 H^1(J,V_\theta)\cong\Hom_B(D,V_\theta).
\]
The vanishing in degree two is what retains all restrictions even
when the source extension is nonsplit.

Let $a_I$ be the multiplicities of represented natural simple types
in the one semisimple head of $D$, and put $A=\sum_Ia_I$.
Their commuting fields are $\F_4$ in the $A$ case and $\F_2$ in
the $S$ case. Each type corresponds to two, respectively six, concrete
top maps. Onto restrictions onto $d$ identical copies are ordered
independent rows over that field. Every cohomology class has exactly
$4^d$ literal cocycles, since the module has no fixed vectors.
The target's first, respectively second, derived group is precisely
its repeated binary module. Thus a map is onto exactly when this
restriction is onto, giving the exact formulas
\begin{align}
 |\Epi(J,R_d^A)|&=2\cdot4^d\sum_I\prod_{j=0}^{d-1}(4^{a_I}-4^j),
 \label{eq:repeated-a}\\
 |\Epi(J,R_d^S)|&=6\cdot4^d\sum_I\prod_{j=0}^{d-1}(2^{a_I}-2^j).
 \label{eq:repeated-s}
\end{align}
A term is zero when $a_I<d$. No independent copy of the source is
assigned to a repeated constituent.

We need the simultaneous constraints
\begin{equation}\label{eq:joint-natural-facets}
 r+2A\leq\lfloor b/2\rfloor\quad(A\text{ case}),\qquad
 r+A\leq\lfloor b/2\rfloor\quad(S\text{ case}).
\end{equation}
Choose bases of all the Hom spaces at once. Their evaluations give
an onto map from $D$ to $V_{\rm all}=\bigoplus_I V_I^{a_I}$.
Lift all its coordinates by the displayed cohomology isomorphism.
The joint image is exactly $V_{\rm all}:R$, where $R$ is the joint
image of the chosen top maps: the restriction of the same $F$ is
onto $V_{\rm all}$, so the full inverse image of $R$ is obtained.
In the $A$ case $R\leq C_3^I$ is odd and acts without trivial
constituents on the binary module. This target has no binary quotient.
Pair its quotient map with the maximal elementary binary quotient
$C_2^r$ of $J$; Goursat's lemma gives the full product. The preimage
of $V_{\rm all}\times C_2^r$ is an actual subgroup of $S_b$ with
an elementary binary quotient of rank $2A+r$. The actual subgroup
abelianization bound \cite[Theorem 2.8]{RoneyDougalTracey2025}
proves the first inequality.

In the $S$ case put $N=R\cap C_3^I$ and $K\cong R/N\leq C_2^I$.
The coordinate sign characters on $N$ are nontrivial and have no
trivial constituent; hence $[N,K]=N$. A complement to the normal odd Hall subgroup gives
$R=N:K$ with elementary binary $K$. The maximal binary quotient
of $V_{\rm all}:R$ is exactly $K$. Pairing with $C_2^r$ therefore
gives $(V_{\rm all}:R)\times C_2^{r-u}$, where $u=\dim K$;
it is the full fibre product over all $u$ binary characters.
Pull back $(V_{\rm all}:K)\times C_2^{r-u}$. Each natural type
restricted to $K$ is a nontrivial transvection module, so its
coinvariant dimension is one. This actual subgroup has binary
abelian quotient rank $A+u+r-u=A+r$, proving the second inequality.

For $q=4$ or $2$,
\[
 \sum_I\prod_{j<d}(q^{a_I}-q^j)
 \leq\sum_I(q^{a_I}-1)^d
 \leq\left(\sum_I(q^{a_I}-1)\right)^d
 \leq(q^A-1)^d<q^{dA}
\]
when the sum is nonempty, and it is zero otherwise.
Above one fixed map to $Q/C$, homomorphic lifts to $Q$ are absent
or a torsor under $\Hom(J,C)$, of size $2^{cr}$. Onto lifts are
only a subset. Consequently the two exponents are bounded by
$cr+2dA$ and $cr+dA$, respectively. Equations
\eqref{eq:joint-natural-facets} bound both by
$\max(c,d)\lfloor b/2\rfloor$. This proves the assertion while
retaining every lift obstruction before the positive enlargement.
\end{proof}

\begin{corollary}[Finite menus with repeated-head entries]
\label{cor:repeated-head-forward}
Fix a finite menu of original transitive actions $U$ of width $w$.
For every original $N\lhd U$, suppose its quotient has either a
fixed mixed character tuple with slope less than $h(w)/8$, or a
specified central repeated-head certificate as in
Proposition~\ref{prop:central-repeat} with
$\max(c,d)/2<h(w)/8$. Then the complete menu has one exponentially
small forward row and zero scalar, with original action weights.
\end{corollary}
\begin{proof}
Fix one choice at every original normal subgroup. The finite set has
a common positive strict gap; sum the estimates of
Propositions~\ref{prop:character-count} and~\ref{prop:central-repeat}
on the same $J$, and apply Lemma~\ref{lem:uniform-delta} to
$\Delta(n,b)Z_J(U)/a(U)$. Different original normals remain different
summands. Repeated-head entries already count literal epimorphisms
and do not receive an additional automorphism factor.
\end{proof}

\begin{lemma}[A binary rank tail on complete sources]
\label{lem:complete-binary-rank-tail}
For each fixed $\epsilon>0$,
\[
 \#\{J\leq S_b:d_2(J)>(1/4+\epsilon)b\}
 \leq2^{b^2/16-\epsilon^2b^2+O_\epsilon(b^{3/2})}.
\]
The same estimate, with an additional $O(b\log(b+2))$ in the
exponent, counts complete $J$ having a normal subgroup $F$ of index
three with $d_2(F)>(1/4+\epsilon)b$.
\end{lemma}
\begin{proof}
For any integer $t\geq0$, the graphs of all maps
$J\to C_2^t$, with the latter on $t$ fixed disjoint pairs, are
subgroups of $S_{b+2t}$. Projection recovers $J$ and its graph.
There are $2^{td_2(J)}$ maps. Apply the coarse bound
\eqref{eq:coarse} and take $t$ within one of $2\epsilon b$.
The exact main exponent is
\[
 \frac{(b+2t)^2}{16}-(1/4+\epsilon)bt
 =\frac{b^2}{16}-\epsilon^2b^2
                       +\frac{(t-2\epsilon b)^2}{4}.
\]
For the second assertion, every index-three overgroup of a fixed
$F\leq S_b$ is generated by $F$ and one permutation, so there are
at most $b!$ possibilities. This is a positive union bound over
all such actual kernels; no eligible kernel multiplicity is divided out.
\end{proof}

\begin{proposition}[All six-pair actions]\label{prop:six-pairs}
Every transitive action on twelve points with six pairs has a complete
exponentially small forward row and a quadratic-exponentially small
scalar, with its original action weight.
\end{proposition}
\begin{proof}
The six-block certificates in Appendix~\ref{sec:certificates} give
$182$ physical action classes and $2476$ literal normal records.
Here is the constructive coverage behind this finite premise.
The $56$ subgroup representatives in $S_6$ contain the identity
and are closed, up to explicit conjugacy, under adjoining every
double-coset representative. If $x=h_1gh_2$, then
$\langle H,x\rangle=\langle H,g\rangle$; induction on generators
therefore proves completeness. Sixteen are transitive. The $2825$
distinct subspaces of $\F_2^6$, the exact Gaussian sum, give all
$112$ invariant-kernel interfaces. On each quotient module,
evaluating a presentation on corrected generator lifts gives the
linear cocycle equations. The number of base-conjugacy classes is
exactly $2^{\dim Z^1-\dim M+\dim M^T}$, proving completeness of
the listed complement graphs and their full preimages, including
nonsplit original lifts. The $224$ transitive block certificates
are identified by actual degree-twelve conjugacy into the $182$
physical actions, with their original symmetric normalizers.

Every original normal list is closed under adjoining the normal
closure of each conjugacy-class representative and contains the
identity. Since a normal subgroup is generated by its conjugacy
classes, this proves exhaustive normal coverage. The certificate
at each normal is either a faithful irreducible tuple consisting
of at most one linear and one general character (two linears
also allowed), or a literal onto map with exactly that kernel to
one of the five fixed models
\[
 C_2^3,\qquad B_{32},\qquad C_2^e\times R_2^A\quad(0\leq e\leq2).
\]
On two four-point blocks, $B_{32}$ is generated by $V_4^2$ and
$(23)(67)$, of order $32$. Kernel intersections of the character
tuples and the explicit model maps verify these alternatives;
labels or group orders alone are not used.

Call the complete $J$ cold if $r=d_2(J)\leq13b/50$.
Character entries cost at most $|\Aut Q|15^{b/3}$.
Every map to a binary target factors through the canonical maximal
binary quotient of $J$. That quotient has generator rank $r$,
so the two binary models have at most $2^{5r}\leq2^{13b/10}$ maps.
Proposition~\ref{prop:central-repeat} bounds each repeated-head
model by $32\cdot2^b$. As $15<16$, summing every original normal
gives a constant $C_U$ with
\[
 Z_J(U)\leq C_U2^{4b/3}\qquad(J\text{ cold}).
\]
This has strict margin $3/2-4/3=1/6$ at original width twelve.
For arbitrary $J$, its generator number is at most $\log_2(b!)$;
the fixed targets have orders at most $46080<2^{16}$. Thus
$Z_J(U)\leq\nu(U)(b!)^{16}$, where $\nu(U)$ is its normal count.
Lemma~\ref{lem:complete-binary-rank-tail} with $\epsilon=1/100$
bounds the hot complete-source count by
$2^{b^2/16-b^2/10000+O(b^{3/2})}$.
The full hot fibre and original normalization are consequently at
most
\[
 \sum_U\frac{\nu(U)(b!)^{16}}{b!a(U)p_nG_{\lfloor n/2\rfloor}}
 2^{b^2/16-b^2/10000+O(b^{3/2})}
 =2^{-\Omega(n^2)},\qquad n=b+12.
\]
The factorials cost only $O(n\log(n+2))$. The cold sum is
$\Delta(n,b)\sum_U C_U2^{4b/3}/a(U)$ and is exponentially small
by Lemma~\ref{lem:uniform-delta}. Fuse the whole family once,
using the existence of any hot accepted pointing and its cold
complement. There is one complete terminal on the cold row.
\end{proof}

\begin{lemma}[A cyclic dual on a common derived source]
\label{lem:binary-cyclic-dual}
Let $X=V:R$, where $V$ is a binary vector space of dimension $d$,
$R$ acts faithfully and $V^*$ is cyclic. Suppose, for $\ell=1$ or
$2$, that $R^{(\ell)}=1$ and the odd part of $R^{(\ell-1)}$ has
zero invariants on $V$. Then
\[
 |\Epi(J,X)|\leq2^d|\GL_d(2)|2^{b/2}.
\]
\end{lemma}
\begin{proof}
The hypotheses give $X^{(\ell)}=V$. Every map therefore sends the
same $F=J^{(\ell)}$ onto $V$. Put $B=J/F$ and $D=H_1(F,\F_2)$.
The odd part of $B^{(\ell-1)}$ maps onto the stipulated odd group
under every onto top map, so averaging gives $H^i(B,V_\theta)=0$
for all $i$. Onto restrictions $D\to V_\theta$ dualize to
injections $V_\theta^*\to D^*$. Fix one cyclic vector in $V^*$.
Its image $f$ determines its whole image module $\F_2[B]f$.
For a fixed $f$, there are at most $|\GL_d(2)|$ linear
identifications; faithfulness determines the top map from each
identification. Thus every top map and onto restriction together
has at most $|\GL_d(2)|2^{\dim D}$ possibilities. Each restriction
has at most $|V|$ cocycle lifts, and the actual $F\leq S_b$ has
$d_2(F)\leq b/2$. This proves the bound for reducible and
nonsemisimple modules as well.
\end{proof}

\begin{proposition}[All four-pair actions]\label{prop:four-pairs}
Every nonbinary transitive action on eight points with four pairs
has a complete exponentially small forward row and a
quadratic-exponentially small scalar, with its original action weight.
\end{proposition}
\begin{proof}
Its top is $A_4$ or $S_4$. In $M=\F_2^4$, the invariant modules
are exactly $0,C,I,M$, with $C$ constant and $I$ augmentation:
weight-two vectors span $I$, weight-one or weight-three orbits
span $M$, and weight four is constant. For each kernel, try all
$256$ pairs of binary corrections to two fixed top generators,
and require that the resulting whole kernel be the selected one.
Every projecting subgroup contains such lifts, so this is exhaustive.
The four-block certificates retain the ten transitive classes,
their original normalizers and all $72$ original normals. Normal
completeness follows from intersection closure of irreducible
character kernels, as in Lemma~\ref{lem:small-pair-data}.
Five classes are
\[
 C_2\times A_4,\quad\operatorname{SL}_2(3),\quad
 S_4\text{ in degree eight},\quad C_2\times S_4,\quad
 \operatorname{GL}_2(3).
\]
The other five, denoted here $Y_1,\ldots,Y_5$, have binary kernels
$I,M,I,I,M$, respectively, with $A_4$ top in the first two cases
and $S_4$ top in the last three. The two augmentation lifts over
$S_4$ are the ordinary lift and the graph whose base parity equals
the top sign; the latter need not split over $I$.

For these five, $C$ is the unique minimal normal subgroup. The
only minimal submodule of their binary kernel is $C$, and any
normal subgroup disjoint from that kernel centralizes it. The
top action on $I$ and $M$ is faithful, so such a subgroup is
trivial. A character nontrivial on $C$, induced and then restricted
to an irreducible constituent, is therefore faithful on $Y_i$.
Thus $|\Epi(J,Y_i)|\leq|\Aut Y_i|5^{b/3}$. The exact central
quotients recorded by the literal maps are
\[
 Y_1/C=R_2^A,\quad Y_2/C=X_c,\quad
 Y_3/C=Y_4/C=X_g,\quad Y_5/C=X_s,
\]
where, on $V=\F_4^2$, $a$ is common multiplication by $\omega$,
$t$ common Frobenius and $s$ coordinate swap, and
\[
 X_c=V:\langle a,s\rangle,\quad
 X_g=V:\langle a,ts\rangle,\quad
 X_s=V:\langle a,t,s\rangle.
\]
Because every nontrivial normal contains $C$, in each case
$Z_J(Y_i)=|\Epi(J,Y_i)|+Z_J(Y_i/C)$ exactly.

For $X_c,X_s$, the module is cyclic over $\F_4[s]$ with unique
proper nonzero submodule the diagonal line. The invariant trace
form $\operatorname{Tr}(x_1y_1^2+x_2y_2^2)$ identifies it with
its binary dual. For $X_g$, the three invariant $\F_4$ lines are
$y=cx$ with $c^3=1$; any two are complementary natural $S_3$
modules. Their commuting field is $\F_2$, and the matrix algebra
$\mathrm{Mat}_2(\F_2)$ shows that two natural copies have a
cyclic dual (take components forming a basis). Thus
Lemma~\ref{lem:binary-cyclic-dual} applies with $\ell=1$ for
$X_c$ and $\ell=2$ for the other two, giving onto bounds
$C_0 2^{b/2}$, $C_0=16|\GL_4(2)|$.
Their exact full normal menus are
\begin{align*}
 Z_J(X_c)&=|\Epi(J,X_c)|+Z_J(C_2\times A_4),\\
 Z_J(X_g)&=|\Epi(J,X_g)|+3|\Epi(J,S_4)|+Z_J(S_3),\\
 Z_J(X_s)&=|\Epi(J,X_s)|+Z_J(C_2\times S_4).
\end{align*}
Indeed in the first and third case every nontrivial normal
contains the diagonal line. In the second case the three line
kernels give three $S_4$ quotients, and all other proper normals
contain $V$. These yield faithful degree-six comparators
$C_2\times A_4,S_4,C_2\times S_4$ (padding $S_4$), with
multipliers $1,3,1$. Since both $1/2$ and $\gamma_5$ are less than
one, Theorem~\ref{thm:finite-comparator} handles $Y_2,\ldots,Y_5$.

For $C_2\times T$, $T=A_4,S_4$, normals not containing the
natural $V_4$ lie in the central $C_2$. Put $z=2^{d_2(J)}$.
The complete remaining normal sum is $z|\Epi(J,A_4)|$ in the
alternating case and $(z-1)|\Epi(J,S_4)|$ in the symmetric case:
in the latter, the central binary row must be independent of the
existing sign. The joint facets in
Proposition~\ref{prop:central-repeat}, with $d=1$, bound these
sums by $8\cdot2^{b/2}$ and $24\cdot2^{b/2}$ respectively.
Their complete quotient intervals are exactly
$Z_J(C_2\times C_3)$ or $Z_J(C_2\times S_3)$, with faithful
five-point representations. This proves the finite comparator
estimate for these two classes. For $\operatorname{SL}_2(3)$
and $\operatorname{GL}_2(3)$ the center $C_2$ is unique minimal
normal, so the same faithful-character argument and exact normal
interval give comparators $A_4,S_4$ of degree four, with slope
$\gamma_5<1$. The nonnatural $S_4$ itself has that comparator with
zero error. This handles all classes except $Y_1$.

For $Y=Y_1$, the repeated quotient has five distinct $\F_4$-line
kernels and therefore
\begin{equation}\label{eq:y-one-full}
 Z_J(Y)=|\Epi(J,Y)|+|\Epi(J,R_2^A)|
                  +5|\Epi(J,A_4)|+Z_J(C_3).
\end{equation}
The first-derived proof of Proposition~\ref{prop:central-repeat}
gives one common head with natural multiplicities $a_I$ and
$r=d_2(J)$. For each corresponding $\theta:J\to C_3$, put
$H=\ker\theta$. Averaging on $J/H=C_3$ gives
\[
 H^1(J,\F_2)=H^1(H,\F_2)^{C_3},\qquad
 H^1(J,V_\theta)=\Hom_{C_3}(H_1(H,\F_2),V_\theta).
\]
Comparing the latter with the same common-head Hom space shows
that $H_1(H,\F_2)$ has trivial multiplicity $r$ and natural
multiplicity $a_I$. Thus
\begin{equation}\label{eq:index-three-head}
 d_2(\ker\theta)=r+2a_I.
\end{equation}
There are at most $b/4$ represented types, since
$2\sum_Ia_I\leq d_2(J')\leq b/2$.

Declare $J$ hot when \emph{some} normal subgroup of index three
has rank $>51b/200$; test all such subgroups, even those whose
natural type is absent. On the cold set, \eqref{eq:index-three-head}
and \eqref{eq:repeated-a} give
\[
 5|\Epi(J,A_4)|\leq10b\,2^{51b/200},\qquad
 |\Epi(J,R_2^A)|\leq8b\,2^{51b/100}.
\]
Together with $Z_J(C_3)\leq3^{b/3}$ and the faithful-character
bound on $Y$, the exact full sum \eqref{eq:y-one-full} has slope
at most $\gamma_5<1$, up to a fixed linear factor in $b$.
For all sources that full sum is at most $C(1+b)2^b$:
use $a_I\leq b/4$ in the repeated-head formula. Lemma
\ref{lem:complete-binary-rank-tail}, now with $\epsilon=1/200$
and its index-three clause, gives the complete hot count
$2^{b^2/16-b^2/40000+O(b^{3/2}+b\log(b+2))}$.
Multiplying by the whole fibre and by the original factor
$1/(384b!p_nG_{\lfloor n/2\rfloor})$, $n=b+8$, leaves a
quadratic-exponentially small scalar. The cold row is exponentially
small by Lemma~\ref{lem:uniform-delta}. The divisor $384$ is the
original $S_8$ normalizer: $Y$ is normal in $C_2\wr S_4$, and its
characteristic $C$ forces every normalizer to preserve those four
pairs. The finite union of all ten classes is fused once, with
all original weights and complete cold terminals retained.
\end{proof}
